\section{Seleção}
O algoritmo \textit{Seleção} também é um método que oferece uma solução simples e intuitiva para o problema da ordenação. Ele funciona particionando o vetor em duas partes: a dos elementos menores em ordem e a dos elementos restantes a serem processados, nesta ordem. A cada iteração, o algoritmo encontra o menor elemento na segunda parte e o troca com o primeiro elemento dela, expandindo a parte ordenada. A implementação em pseudo-código abaixo garante ordenação estável.

\lstinputlisting[language=C]{funcoes-c/d_selecao.txt}
% \begin{algorithm}
%     \caption{Seleção}\label{alg-selection-sort:cap}
%     \DontPrintSemicolon
%     \Entrada{um vetor $v$ de $n$ elementos}
%     \Resultado{o vetor $v$ em ordem crescente}
%     \Inicio{
%         \Para{$i = 0$ \Ate $\,n - 2\,$}{
%             $indice\_do\_menor\_elementor \gets i$\;
%             \Para{$j = i+1$ \Ate $\,n - 1\,$}{
%                 \Se{$v[j] < v[indice\_do\_menor\_elementor]$}{
%                     $indice\_do\_menor\_elementor \gets j$\;
%                 }
%             }
%             Troca($v[i]$, $v[indice\_do\_menor\_elementor]$)\;
%         }
%     }
% \end{algorithm}

\subsection{Prova de correção}
No início de cada iteração do laço externo, vale o seguinte invariante: o vetor $v[0..i-1]$ contém os menores elementos e em ordem. Este invariante será usado para provar, por indução em $i$, que o algoritmo é correto.

Para o caso base, ou seja, no início da primeira iteração, quando $i = 0$, o vetor $v[0..i-1]$ é vazio e atende o invariante por definição.

Agora suponha que no início da $(i+1)-$ésima iteração o invariante seja válido, isto é, o vetor $v[0..i-1]$ contém os $i$ menos elementos do vetor original e em ordem. Então observe que o laço interno vai encontrar o índice do menor elemento de $v[i..n-1]$ e, por fim, este elemento será trocado com o elemento da posição $i$. Logo, ao final da iteração, $v[0..i]$ estará ordenado e o invariante permanece válido para a próxima iteração.

Portanto, pelo princípio da indução, o invariante permanece válido em todas as iterações do laço externo. Sendo assim, no início da última iteração, $v[0..n-1]$ estará em ordem, logo o algoritmo é correto.

\subsection{Complexidade de tempo e espaço}
O ponto crítico do \textit{Seleção} é o teste condicional \textbf{se} dentro do laço interno. Assim como no algoritmo \textit{Bolha}, a quantidade de vezes que esse ponto é executado também é dada pela soma \ref{eq:1}. Portanto, este algoritmo sempre requer tempo proporcional a $O(n^2)$, enquanto que o uso de memória adicional é $O(1)$.